\documentclass[11pt]{article}
\usepackage[a4paper,margin=25mm]{geometry}
\usepackage{amsmath,amssymb,amsthm,mathtools}
\usepackage[T1]{fontenc}
\usepackage{hyperref}
\hypersetup{colorlinks=true,linkcolor=blue,urlcolor=blue}
\newtheorem{theorem}{Theorem}
\newtheorem{lemma}[theorem]{Lemma}
\newtheorem{corollary}[theorem]{Corollary}
\newcommand{\Tr}{\operatorname{Tr}}
\newcommand{\EB}{\operatorname{EB}}
\newcommand{\id}{\operatorname{id}}
\newcommand{\norm}[1]{\left\lVert#1\right\rVert}
\title{Linear classicalization of approximately clonable quantum experiments\\and sharp transport stability}
\author{ASTRA ULTRA research investigation}
\date{11 October 2026}
\begin{document}
\maketitle

\noindent\textbf{Conclusion and status.}
The leading completed deduction classically reconstructs an approximately product-clonable quantum experiment, including mixed states, with dimension-independent linear infidelity. One fixed rank-one measurement with at most $2d$ outcomes works simultaneously on every input state for the supplied cloner. A quadratic risk gap, deletion induction, constrained decision amplification, and spectral polar rounding perform the necessary work. For pure families the sharper outcome budget is $d+O(d\varepsilon)$, with a matching-order lower construction. A second deduction gives sharp stability of the Brenier-map embedding, including exponent $1/3$ for a full Gaussian source with just one higher-moment target. Complete written arguments, exact counterexamples, supplied interfaces, and source dependencies are preserved here and in the companion proof files.

The same investigation gives a geometric classification of positive finite-memory composition: the minimum nonzero orbit dimension determines the sharp lifetime exponent. It yields an exact classical memory law for a fixed normalized-trace experiment, with one preparation and one binary measurement, and an optimal qutrit realization of that process.

These are internally reconstructed mathematical results. Historical originality, external validation, and an independently major foundational consequence have not been established. No result is presented as the requested historic or foundational breakthrough. In particular, mathematical existence does not supply an unknown-state acquisition procedure or a dimension-independent algorithm.

\tableofcontents

\section{Pinned capital and the distinction between sources and deductions}

The read-only source revisions are
\begin{itemize}
\item Astra: \texttt{b2a7a8fa91a8fcf6da05f36100ae4ce64a72a729};
\item OpenAI Math: \texttt{fd4aeeb2ee4fc729c18d98444fed42fd0529eeeb}.
\end{itemize}
The manifest also pins the unmerged Astra pull-request heads 41--47. Astra's routing files, current status, dossiers, proof ranges, and selected original proof assets were read. OpenAI's overview, manuscript catalogue, original sources, available formalization routes, reasoning material, and revision history supplied breadth; selected complete proof mechanisms supplied depth. This was a selective proof investigation, not an assertion that every manuscript was read or certified. The October revision and withdrawal notices were treated as material scientific status.

The Astra predecessor for the original pure-state route in Section~\ref{sec:clone} is N407--N408, particularly the finite Schur-partition proof, preserved at
\href{https://github.com/harjassand/astra-research-knowledge/blob/b2a7a8fa91a8fcf6da05f36100ae4ce64a72a729/web/pages/d-2115f1527533d2f7.html}{the pinned N407 original source route}.
Its original indexed SHA256 is
\begin{center}\small\texttt{984acfe2c372aea82445a8bb4b4c9034dd0671469d5615e8f5b78838f8371e18}.
\end{center}
That proof's scientific status remains source-derived and unreviewed in the repository. We reconstruct its finite argument below; repository inclusion is not a premise of correctness. N410's cube-root average bound and N534's uniform linear reconstruction are comparators, not black-box premises for the new average bound. The later mixed-state theorem in Section~\ref{sec:mixed} instead uses a direct risk-gap/deletion argument and does not require the Schur partition theorem or a purification lift.

The transport predecessor is OpenAI family 374,
\href{https://github.com/openai/math/tree/fd4aeeb2ee4fc729c18d98444fed42fd0529eeeb/preprints/Sharp-One-Third-Stability-of-Brenier-Maps-September-25-2026/build/source/sections}{Sharp One-Third Stability of Brenier Maps}, especially \texttt{cells.tex}, \texttt{moments.tex}, \texttt{finite\_stability.tex}, \texttt{gradient.tex}, and \texttt{sharpness.tex}. Its compact uniform-source endpoint is credited to that source. The new deductions are weighted midpoint domination, interpolation toward slice modes, the one-sided tail transfer, and a sharp general radial-budget modulus.

\section{A dimension-free mixed-state classicalization theorem}\label{sec:mixed}

Write $F(\rho,\sigma)=\norm{\sqrt\rho\sqrt\sigma}_1^2$ for squared Uhlmann fidelity. The premise below concerns closeness to an independent \emph{joint product}, rather than merely accurate output marginals.

\begin{theorem}[One classical surrogate for every input state]\label{thm:mixed}
Let $H$ have finite dimension $d$, and let $B:H\to H\otimes H$ be any channel. There exist an isometry $U:H\to\mathbb C^N$, $d\le N\le2d$, and fixed density matrices $\sigma_1,\ldots,\sigma_N$ such that the entanglement-breaking channel
\[
\Lambda_B(\rho)=\sum_{j=1}^N\langle j|U\rho U^*|j\rangle\sigma_j
\]
satisfies, simultaneously for every density matrix $\rho$,
\[
1-F(\rho,\Lambda_B(\rho))\le2800\,[1-F(B(\rho),\rho\otimes\rho)].
\tag{M1}
\]
Thus one ancillary qubit suffices for a coordinate-projective dilation of the \emph{measurement}. Neither acquisition of $U$ and the preparations nor their circuit cost is supplied for free.
\end{theorem}

The following proof supersedes the earlier mixed-state gap recorded in the exploratory notes. No purification of the input family is assumed. Put $e(\rho)=1-F(B(\rho),\rho^{\otimes2})$, $\tau=1/2800$, and $S_\tau=\{\rho:e(\rho)\le\tau\}$.

\subsection{Product cloning forces a separated class geometry}

For $\rho,\sigma\in S_\tau$, put $r=\sqrt{F(\rho,\sigma)}$. Contractivity and the triangle inequality for the Bures angle give
\[
g(r):=\arccos(r^2)-\arccos r
\le\arcsin\sqrt{e(\rho)}+\arcsin\sqrt{e(\sigma)}.
\tag{M2}
\]
Since $g(r)\ge r(1-r)$ and $\arcsin x<(6/5)x$ in the admitted range, $r(1-r)<1/20$. Hence either $r<1/16$ or $r>15/16$. The high-overlap relation is transitive: two high overlaps give a third greater than $\cos(2\arccos(15/16))=97/128>1/16$, forcing it high. It partitions the compact good set into finitely many compact classes. Finiteness follows from uniform Bures separation in a compact finite-dimensional state space.

Within a class, $r>3/4$, and
\[
g(r)\ge\frac{r(1-r)}{\sqrt{1-r^4}}
=\frac{r}{(1+r)\sqrt{1+r^2}}\sqrt{1-r^2}
\ge\frac{12}{35}\sqrt{1-r^2}.
\]
Combining with (M2) gives $1-F(\rho,\sigma)\le25(e(\rho)+e(\sigma))$. Each class has a minimum-error representative $\rho_c$, and consequently
\[
1-F(\rho,\rho_c)\le50e(\rho),\qquad e(\rho_c)\le e(\rho).
\tag{M3}
\]

\subsection{A quadratic risk gap and deletion induction}

Take $m$ representatives from distinct classes. Let $R_m$ be their optimal worst-case discrimination error. A minimizer exists in the compact POVM set. Two uses of $B$ produce three outputs within Bures distance $\sqrt{8e(\rho)}$ of $\rho^{\otimes3}$. For any effect, its square-root probability changes by at most the Bures distance. Measuring the minimizing POVM on the three ideal copies and taking majority has error at most $3R_m^2$. Pulling this decision back through the actual two-call channel proves
\[
R_m\le f(R_m),\qquad f(x)=(\sqrt3\,x+\sqrt{1/350})^2.
\tag{M4}
\]
Set $a=1/200$ and $b=13/50$. The convex quadratic $f(x)-x$ is negative throughout $[a,b]$. Explicitly, $\sqrt3<7/4$ and $1/\sqrt{350}<3/56$ give
\[
f(a)<(349/5600)^2<1/256<a,\qquad
f(b)<(89/175)^2<b.
\]
Thus any attainable upper bound $R_m\le b$ forces $R_m<a$.

For $m\le3$ the pretty-good measurement has prior-average error at most
\[
\sum_{i\ne j}\sqrt{p_ip_j}\sqrt{F(\rho_i,\rho_j)}
\le(m-1)/16\le1/8.
\tag{M5}
\]
This bound holds for every prior, so finite minimax gives $R_m\le1/8<b$. Here (M5) can be proved directly: put $A_i=p_i\rho_i$, let the positive block Gram matrix have blocks $G_{ij}=\sqrt{A_i}\sqrt{A_j}$, and let $G_0=\operatorname{diag}(A_i)$. The PGM error is $\sum_{i\ne j}\norm{(\sqrt G)_{ij}}_{\rm HS}^2$, bounded by $\norm{\sqrt G-\sqrt{G_0}}_{\rm HS}^2\le\norm{G-G_0}_1\le\sum_{i\ne j}\norm{G_{ij}}_1$. The middle inequality is Powers--St\o rmer; the bound $(\sum_i\sqrt{p_i})^2-1\le m-1$ completes (M5). Singular states and zero priors are handled on the support, with arbitrary complementary effects.

Inductively, for $m\ge4$, randomly omit one label and use the already constructed $(m-1)$-label measurement. Each input is omitted with probability $1/m$, so
\[
R_m\le1/m+(1-1/m)a\le203/800<b.
\tag{M6}
\]
Equation (M4) forces $R_m<a$. This proves a dimension-independent weak classifier for any finite number of representatives. It also bounds the total number of classes: success $>1-a$ on each representative gives $k(1-a)<\sum_c\Tr M_c=d$, hence $k<200d/199$.

The representative measurement extends to its whole class by (M3) and the fidelity/trace-distance bound:
\[
r(\rho)\le a+\sqrt{50\tau}=1/200+1/\sqrt{56}<1/4<b.
\]
Minimize the worst error over the entire compact good set. The same majority argument (M4) then forces its risk below $a$. We have therefore constructed one class POVM with error at most $a$ on every good mixed state.

\subsection{Constrained amplification turns the gap into linear loss}

Keep the compact convex set of class POVMs whose errors on every good state are at most $a$. For any finite prior on that set, minimize average class error $R$. Let $E$ be its average cloning infidelity. Majority preserves all constraints because $f(a)<a$, and the pointwise pulled-back error obeys
\[
r_i'\le(\sqrt3\,r_i+\sqrt{8e_i})^2.
\]
Optimality, $r_i^2\le ar_i$, and Cauchy--Schwarz imply
\[
R\le(\sqrt{3aR}+\sqrt{8E})^2,
\qquad R\le C_*E,\quad C_*=\frac8{(1-\sqrt{3/200})^2}<11.
\tag{M7}
\]
Finite minimax and the finite-intersection property within the same globally constrained POVM set give one class measurement with $r(\rho)\le C_*e(\rho)$ simultaneously on every good state. Prepare $\rho_c$ on class $c$. Separate concavity of squared fidelity and (M3) give reconstruction loss at most $(C_*+50)e(\rho)<61e(\rho)$. Bad states have loss at most $1<2800e(\rho)$. This already proves a fixed global EB surrogate; the next operation gives the stated measurement size.

\subsection{Spectral polar rounding uses at most twice the input dimension}

For the class POVM just obtained, set $P_c=\mathbf1_{[1/2,1]}(M_c)$. Then $P_c\le2M_c$, $I-P_c\le2(I-M_c)$, and $S=\sum_cP_c\le2I$. Let $T:H\to\bigoplus_c\operatorname{ran}P_c$ stack the projections, and pad its output to dimension $N=\max(d,\sum_c\operatorname{rank}P_c)\le2d$. Complete the polar factor to an isometry $U$, so $T=U\sqrt S$.

For any input $\rho$ in class $c$, write $r=1-\Tr(M_c\rho)$, $A=\sum_{j\ne c}P_j$, and $Q=I-P_c$. Then $\Tr(\rho A)\le2r$, $\Tr(\rho Q)\le2r$, and
\[
(S-I)^2=(A-Q)^2\le2A^2+2Q^2\le4A+2Q.
\]
Consequently $\norm{(U-T)\sqrt\rho}_{\rm HS}^2=\Tr[\rho(I-\sqrt S)^2]\le12r$, while the wrong-block energy of $T\sqrt\rho$ is at most $2r$. Their triangle inequality proves wrong-block probability at most $(14+4\sqrt6)r<24r$ after $U$.

Refine the output blocks into individual coordinate rays, preparing $\rho_c$ on every coordinate in block $c$, and any representative on padded coordinates. All nonzero effects have rank one, and there are at most $2d$ outcomes. The good-state reconstruction loss is below $(24C_*+50)e(\rho)<314e(\rho)$; bad states are covered by $2800e(\rho)$. This proves (M1). Zero coordinate effects may be discarded; if the good set is empty, an ordinary basis measurement with a fixed output state suffices.

\paragraph{Resource and consequence boundary.}
The conclusion supplies one channel for every input and every prior, including exactly clonable mixed states. For any finite mixed ensemble it yields optimal reconstruction infidelity at most $2800$ times optimal product-cloning infidelity. Its measurement dilation fits $H\otimes\mathbb C^2$; mixed-output preparation may use classical randomness or additional state-preparation resources, which are separate and unspecified. The global class representatives and constrained POVM have unknown acquisition cost. No black-box tomography, efficient circuit, diamond-norm simulation, or channel-capacity theorem follows. The complete root and independent derivations are retained in \texttt{proofs/mixed\_classicalization.txt} and \texttt{proofs/mixed\_independent.txt}.

\subsection{A commuting core after one small dilation}

The same isometry gives a stronger representation statement. Let $\Pi_c$ project onto output block $c$, and define
\[
p_c=\Tr(\Pi_cU\rho_cU^*),\qquad
\Omega_c=\Pi_cU\rho_cU^*\Pi_c/p_c.
\]
All $p_c>0$. The states $\Omega_c$ have mutually orthogonal supports and therefore commute exactly. Projecting a purification shows $F(U\rho_cU^*,\Omega_c)\ge p_c\ge1-24C_*e(\rho_c)$. The triangle inequality for purified distance $\sqrt{1-F}$, (M3), and $e(\rho_c)\le e(\rho)$ give, for every good state in class $c$,
\[
1-F(U\rho U^*,\Omega_c)
\le(\sqrt{50}+\sqrt{24C_*})^2e(\rho)<550e(\rho).
\tag{M8}
\]
Thus approximate product replicability forces proximity to one fixed finite dictionary of orthogonal mixed states in a space of dimension at most $2d$. The dictionary and embedding are independent of the input prior. The dictionary states need not lie in $U(H)$; this is a dilation representation, distinct from the output states $\rho_c$ of the physical reconstruction channel on $H$.

\subsection{Marginal broadcasting has no dimension-free reconstruction modulus}

The independent-product premise in (M1) cannot be replaced by accurate marginals, even with an arbitrary modulus tending to zero. Here is a complete separation, using a reservoir of different copy counts. Fix $L\ge2$, $K=2^L$, and $b=K^2$. On
\[
H_L=\bigoplus_{j=0}^L\operatorname{Sym}^{2^j}(\mathbb C^b),\qquad
D_L=\sum_{j=0}^L\binom{b+2^j-1}{2^j},
\]
consider, for every unit vector $\psi\in\mathbb C^b$,
\[
\rho_\psi=\frac1L\bigoplus_{j=1}^LP_\psi^{\otimes2^j}.
\]
A channel measures the orthogonal copy-count block $j$ and splits its existing $2^j$ copies into two groups of $2^{j-1}$ copies, placing both outputs in block $j-1$. Define its action on block zero arbitrarily to complete a channel. Each output marginal is
$\widetilde\rho_\psi=L^{-1}\bigoplus_{j=0}^{L-1}P_\psi^{\otimes2^j}$, so
\[
F(\rho_\psi,\widetilde\rho_\psi)=((L-1)/L)^2,
\qquad e_{\rm marginal}\le2/L\longrightarrow0.
\tag{M9}
\]

Every EB reconstruction channel nevertheless has Haar-average fidelity at most $4/b$, tending to zero. To prove this, let $D_n=\dim\operatorname{Sym}^n(\mathbb C^b)$. For $n$ input copies and any prepared state in a $q$-copy output block,
\[
\int P_\psi^{\otimes(n+q)}d\psi=\Pi_{n+q}/D_{n+q}\le I/D_{n+q}.
\]
For a POVM effect $M_y$ and output block $\sigma_{y,q}$, positivity bounds the integrated support overlap by $\Tr M_y\Tr\sigma_{y,q}/D_{n+q}$. Sum over outcomes and all target blocks, using $\sum_y M_y=I$ and $\sum_q\Tr\sigma_{y,q}\le1$. Since $n\le K$ and every target block has $q\ge2$, the total is at most
\[
\max_{n\le K,q\ge2}\frac{D_n}{D_{n+q}}
=\frac{D_K}{D_{K+2}}
=\frac{(K+1)(K+2)}{(K+b)(K+b+1)}\le4/b.
\]
The argument allows reading the input clock and choosing arbitrary output clocks, and averaging over the input mixture preserves the bound. The target $\rho_\psi$ is supported on $A_\psi=\bigoplus_{j=1}^LP_\psi^{\otimes2^j}$. Fidelity monotonicity under its support test gives $F(\rho_\psi,\Lambda(\rho_\psi))\le\Tr[A_\psi\Lambda(\rho_\psi)]$. Therefore
\[
\inf_{\Lambda\in\EB}\int[1-F(\rho_\psi,\Lambda(\rho_\psi))]d\psi
\ge1-4/b\longrightarrow1.
\tag{M10}
\]
This contradicts every proposed dimension-free modulus from marginal broadcasting error to EB reconstruction error. The actual joint output has a perfectly shared clock and product-target fidelity $(L-1)^2/L^3\to0$, so it does not satisfy (M1)'s premise. The dimension and supplied copy reservoir are enormous and explicitly displayed; there is no efficient unknown-state acquisition claim. Haar averaging can be replaced by a finite weighted exact $2K$-design, whose existence follows from finite-dimensional convexity of its moment matrix; no inexpensive design acquisition is asserted. The full independent proof is in \texttt{proofs/marginal\_separation.txt}.

\subsection{Mixed states can retain measurement overhead as their defect vanishes}

The pure-family $d+O(d\varepsilon)$ measurement-space bound below does not extend to mixed states. Fix any finite loss constant $K_0$, choose integer $q>2K_0+1$, and take $b_0$ orthogonal $q$-dimensional simplex blocks with $q+1$ labels each. Attach one orthogonal tag $C_a$ per label, and put $\rho_a(w)=(1-w)C_a+wP_a$. The input dimension is $D=b_0(2q+1)$, with $n=b_0(q+1)$ labels. The POVM $E_a=C_a+\frac q{q+1}P_a$ sums to $I$. Measuring it and preparing $\rho_a(w)^{\otimes2}$ gives product-cloning error at most $\varepsilon_w=w/(q+1)\to0$.

For all sufficiently small $w>0$, every surrogate with loss at most $K_0\varepsilon_w$ on every label needs a fixed projective Naimark measurement space of dimension at least $D+b_0$. To see this, take a hypothetical smaller-dimensional sequence as $w\to0$. Compactness gives a limiting measurement/preparation channel exactly fixing all pure tags. Grouping outcomes preparing the same tag shows that these tags occupy $n$ orthogonal measurement directions; the remaining minority inputs have at most $N-n$ effective preparations. The target support test $\Pi_a=C_a+P_a$ gives, after dividing the loss inequality by $w$ and taking the limit,
$\Tr[\Pi_a\Lambda_0(P_a)]\ge1-K_0/(q+1)$. Since $\norm{\Pi_a+\Pi_j}_\infty\le1+1/q$, one preparation cannot serve two labels at this threshold. Hence $N-n\ge n$, or $N\ge D+b_0$, a contradiction. The full compactness and residual-projection argument is in \texttt{proofs/mixed\_overhead.txt}. The small-$w$ cutoff is existential. This concerns one fixed isometry/projective measurement, and allows arbitrary mixed preparations; externally randomized choices among measurements require a separate resource model.

\section{A sharper pure-state frame theorem}\label{sec:clone}

All Hilbert spaces in this section are finite-dimensional. A channel is completely positive and trace preserving. For a pure state $P_\psi=|\psi\rangle\langle\psi|$, ``infidelity'' means $1-\Tr[P_\psi\rho]$; it is the loss associated with squared fidelity against a pure target.

\begin{theorem}[Universal linear classical surrogate]\label{thm:surrogate}
For every channel $B:H\to H\otimes H$, there is one entanglement-breaking channel $\Lambda_B:H\to H$ such that, simultaneously for every unit vector $\psi\in H$,
\[
1-\Tr[P_\psi\Lambda_B(P_\psi)]
\le 2800\bigl(1-\Tr[(P_\psi\otimes P_\psi)B(P_\psi)]\bigr).
\tag{1}
\]
Consequently, for every finite pure ensemble $\mathcal E=\{(p_i,P_i)\}$, define
\begin{align*}
e_{\rm cl}(\mathcal E)&=1-\max_B\sum_i p_i\Tr[(P_i\otimes P_i)B(P_i)],\\
e_{\EB}(\mathcal E)&=1-\max_{\Lambda\in\EB}\sum_i p_i\Tr[P_i\Lambda(P_i)].
\end{align*}
Then
\[
e_{\EB}\le \min\{1,2800 e_{\rm cl}\},\qquad e_{\rm cl}\le 2e_{\EB}.
\tag{2}
\]
The constants do not depend on dimension, number of labels, or smallest positive prior.
\end{theorem}

We prove the finite-ensemble inequality for a fixed channel, then obtain (1) by minimax. Put $e_i=1-\Tr[(P_i\otimes P_i)B(P_i)]$ and $\tau=1/2800$.

\subsection{A weak class measurement from approximate replication}

First suppose $e_i\le\tau$ for every label under consideration. A Stinespring isometry $V:H\to H\otimes H\otimes E$ gives unit vectors
\[
u_i=V\psi_i,\qquad v_i=\psi_i\otimes\psi_i\otimes f_i,
\qquad \norm{u_i-v_i}^2=2(1-\sqrt{1-e_i})\le2e_i.
\]
Here $f_i$ is the normalized projection of $u_i$ onto the whole target subspace $\psi_i\otimes\psi_i\otimes E$; its phase makes $\langle v_i,u_i\rangle\ge0$.
Let $C_{ij}=\langle\psi_i,\psi_j\rangle$, $F_{ij}=\langle f_i,f_j\rangle$, and $L=C\circ F$. Schur multiplication by a correlation matrix is unital completely positive and has completely bounded norm one. The label-diagonal Stinespring isometries
\[
W_u|i\rangle=|i\rangle u_i,\qquad W_v|i\rangle=|i\rangle v_i
\]
give, by telescoping at every amplification,
\[
\norm{S_C-S_{C\circ C\circ F}}_{\rm cb}\le2\sup_i\norm{u_i-v_i}
\le\delta:=2\sqrt{2\tau}.
\]
Composition with $S_F$ implies $\norm{S_L^2-S_L}_{\rm cb}\le\delta$.

In the unital Banach algebra of Schur multipliers set
\[
Z=2S_L-I,\quad D=4(S_L^2-S_L),\quad
R=(I+D)^{-1/2},\quad S_Q=(I+ZR)/2.
\]
The binomial series converges since $4\delta<1$. As $Z^2=I+D$, the factors commute and $S_Q^2=S_Q$. The symbol $Q$ is Hermitian, has diagonal one, and is entrywise zero or one. Also
\[
\norm{S_Q-S_L}_{\rm cb}\le
\eta:=\frac32\bigl((1-4\delta)^{-1/2}-1\bigr).
\]
Exact sufficient bounds are $\delta<27/500$, $\eta<49/250$, so $a:=\delta+\eta<1/4$. For example, the second follows from
$1-4(27/500)=98/125$ and $\sqrt{125/98}<424/375$.

The ones of $Q$ form an equivalence relation. Otherwise a principal three-label mask is
\[
\begin{pmatrix}1&1&0\\1&1&1\\0&1&1\end{pmatrix}.
\]
Applied to the all-ones rank-one density matrix divided by three, its output trace norm is $(1+2\sqrt2)/3>5/4$. Trace duality would give $\norm{S_Q}_{\rm cb}>5/4$, contrary to $\norm{S_Q}_{\rm cb}\le1+\eta<5/4$. Thus $Q$ specifies fixed classes. Within each class $|C_{ij}|\ge1-\eta>3/4$.

Moreover
\[
\norm{S_C-S_{C\circ Q}}_{\rm cb}\le a.
\]
For any prior $s$ on this fixed family, let $G=\operatorname{diag}(\sqrt s)C\operatorname{diag}(\sqrt s)$ and $G_0=G\circ Q$. Trace duality gives $\norm{G-G_0}_1\le a$; complex conjugation handles the transpose in the Schur predual. Powers--St\o rmer gives
\[
\norm{\sqrt G-\sqrt{G_0}}_2^2\le a.
\]
The vectors $\alpha_i=\sqrt G|i\rangle/\sqrt{s_i}$ have Gram matrix $C$, whereas $\beta_i=\sqrt{G_0}|i\rangle/\sqrt{s_i}$ are unit vectors in their own coordinate classes. Pulling the class projections back along the isometry $\psi_i\mapsto\alpha_i$ gives a class POVM with prior-average wrong-class probability at most $a$. Zero prior labels are handled by restriction and extension, or by positive-prior approximation. Since this holds for every prior on the same classes, finite minimax gives one full-domain class POVM with
\[
r_i\le a_0:=1/4\quad\hbox{for every good label.}\tag{3}
\]
The reconstructed source mechanism stops at a weak class measurement. The next two steps retain the individual cloning losses.

\subsection{Within-class loss is controlled by individual errors}

For two good states put $r=|\langle\psi_i,\psi_j\rangle|>3/4$, $s=\sqrt{1-r^2}$, and $E_{ij}=e_i+e_j$. Decomposing $u_i=\sqrt{1-e_i}\,v_i+w_i$, with $w_i$ orthogonal to the entire corresponding target subspace, and using preservation of the input overlap gives
\[
r\le r^2+(\sqrt{e_i}+\sqrt{e_j})\sqrt{1-r^4}+\sqrt{e_i e_j}.
\]
The target--residual terms have the factor $\sqrt{1-r^4}$ precisely because of that full-subspace orthogonality. Hence
\[
\frac37s^2\le2\sqrt{E_{ij}}s+E_{ij}/2.
\]
If $E_{ij}>0$, $x=s/\sqrt{E_{ij}}\ge5$ would imply
$15/7\le3x/7\le2+1/(2x)\le21/10$, a contradiction. The zero-error case follows directly. Therefore
\[
1-|\langle\psi_i,\psi_j\rangle|^2\le25(e_i+e_j).\tag{4}
\]
Choose in each class a representative with minimal $e_i$. Its preparation loss on input $i$ is at most $50e_i$.

\subsection{Constrained decision amplification closes the average bound}

For a conditional prior on the good labels, minimize $R=\sum_i p_i r_i$ over class POVMs subject to every constraint $r_i\le1/4$. Compactness and (3) give an optimizer. Apply $B$ twice, splitting one of its outputs, to obtain a physical one-to-three channel. At Stinespring level,
\[
\norm{U_3\psi_i-\psi_i^{\otimes3}\otimes f_i\otimes f_i}
\le2\sqrt{2e_i}.\tag{5}
\]
This follows by two applications of the isometric triangle inequality.

Measure the optimizing class POVM on all three outputs. Return a class occurring twice, using an arbitrary fixed rule if all three labels differ. On the ideal product target, an incorrect decision requires at least two individual errors, so it has probability at most $3r_i^2$. Applying the square root of the wrong-decision effect to (5) yields a POVM on the original input with error
\[
r_i'\le(\sqrt3\,r_i+\sqrt{8e_i})^2.\tag{6}
\]
This new POVM remains feasible: $\sqrt3/4+\sqrt{8\tau}<7/16+1/16=1/2$. Independence was used only on the ideal product state; no independence of the actual cloned outputs is assumed.

Optimality, $r_i^2\le a_0r_i$, and Cauchy--Schwarz give, with $E=\sum_i p_i e_i$,
\[
R\le\sum_i p_i r_i'
\le\bigl(\sqrt{3a_0R}+\sqrt{8E}\bigr)^2.
\]
Taking square roots and rearranging proves
\[
R\le\frac8{(1-\sqrt3/2)^2}E=(224+128\sqrt3)E.\tag{7}
\]
Preparing the class representatives now costs at most
\[
(274+128\sqrt3)E<500E.
\]
For the original ensemble, call $e_i>\tau$ bad. Their total prior mass is at most $2800\sum_{\rm bad}p_i e_i$. The same full-domain EB channel therefore has total loss at most
\[
500\sum_{\rm good}p_i e_i+2800\sum_{\rm bad}p_i e_i
\le2800\sum_i p_i e_i.
\]
If there are no good labels the bound is trivial. This proves the finite-ensemble inequality, including exact zero error.

For a fixed finite family and fixed $B$, apply minimax to the linear losses
$1-\Tr[P_i\Lambda(P_i)]-2800e_i$ on the compact convex EB-channel set. One $\Lambda$ makes every loss nonpositive. The finite-intersection property on the same compact set extends this to all pure states, proving (1). Finally, an EB channel with output states $\sigma_y$ can prepare $\sigma_y\otimes\sigma_y$ instead; $1-f^2\le2(1-f)$ proves the other inequality in (2).

\subsection{The rate is optimal and a Gram gate is closed}

For two equally likely pure states with overlap $c$, the optimal measure-and-prepare fidelity is
\[
F_{\EB}=\frac{1+\sqrt{1-c^2+c^4}}2.
\]
This formula is established in \href{https://arxiv.org/abs/quant-ph/0302092}{Fuchs and Sasaki, accessible fidelity}. An explicit isometry aligning the input pair of projective angle $\theta=\arccos c$ with the ideal two-copy pair of angle $\phi=\arccos c^2$ attains cloning infidelity $\sin^2((\phi-\theta)/2)$. Thus, as $\theta\downarrow0$,
\[
e_{\EB}=\theta^2/4+o(\theta^2),\qquad
e_{\rm cl}\le(3-2\sqrt2)\theta^2/4+o(\theta^2).
\]
Any universal constant in (2) is at least $3+2\sqrt2$. Linear order cannot be improved to $o(e_{\rm cl})$ uniformly. The optimal constant remains unresolved.

\subsection{One surrogate needs only linearly many pure preparations}

There is a stronger structural conclusion: $\Lambda_B$ in (1) can have at most $\lfloor8d/7\rfloor$ outcomes with pure prepared states, where $d=\dim H$. If no pure input has $e_B(P)\le\tau$, one outcome suffices. Otherwise the good projective set is compact. The quantity $L(P,Q)=C(P,Q)F(P,Q)$ above is invariant under vector phases and continuous on the good set, since its normalization denominators are at least $1-\tau$. Its values lie within $\eta$ of zero or one. Thus $\operatorname{Re}L>1/2$ gives a global equivalence relation whose classes are closed.

Choose any $k$ distinct class representatives. Their uniform weighted Gram matrix has rank at most $d$, and its block pinching is $I_k/k$, with $\norm{G-I_k/k}_1\le a<1/4$. Testing on the kernel of $G$, and using that the difference has trace zero, gives
\[
2(k-d)/k\le\norm{G-I_k/k}_1<1/4.
\]
Hence there are at most $\lfloor8d/7\rfloor$ classes. Each is compact and has a minimum-error representative. A fixed finite-outcome class-POVM set satisfying $r(P)\le1/4$ for every good $P$ is compact and nonempty by the finite-intersection property. Optimize a finite prior over this global feasible set; majority preserves every constraint. The same argument as (7), then finite minimax and compactness, gives one class POVM with $r(P)\le(224+128\sqrt3)e_B(P)$ simultaneously. Preparing the minimum-error representatives proves the claimed outcome bound and (1). This provides a finite representation; finding the global classes or the global constrained optimum is not supplied as an efficient algorithm.

\subsection{Pointwise polar compression gives a Parseval frame}

Write $C_{\rm dec}=224+128\sqrt3<450$, and take the global class POVM $\{M_c\}_{c=1}^k$ and minimum-error representatives $P_c=|\psi_c\rangle\langle\psi_c|$ just constructed. Put
\[
s_c=\langle\psi_c,M_c\psi_c\rangle,\quad
m_c=M_c\psi_c/\sqrt{s_c},\quad N=\max(d,k).
\]
All $s_c$ are positive. Since $M_c^{1/2}\psi_c/\sqrt{s_c}$ is unit,
\[
|m_c\rangle\langle m_c|\le M_c,\quad
\sum_c|m_c\rangle\langle m_c|\le I,\quad
|\langle m_c,\psi_c\rangle|^2=s_c.
\]
Let $T:H\to\mathbb C^N$ have rows $m_c^*$ and zero padded rows. Set $S=T^*T\le I$. Complete its polar partial isometry to an isometry $U$, possible because $N\ge d$, so $T=U\sqrt S$. For $r_c=1-s_c$,
\[
\norm{(U-T)\psi_c}^2
=\langle\psi_c,(I-\sqrt S)^2\psi_c\rangle
\le1-\norm{T\psi_c}^2\le r_c.
\]
The squared norm of the wrong coordinates of $T\psi_c$ is also at most $r_c$. Their triangle inequality gives
\[
1-|\langle c,U\psi_c\rangle|^2\le4r_c.
\]
For any good $P\in S_c$, the pure-state trace-distance triangle, the within-class bound, and $e_B(P_c)\le e_B(P)$ therefore give
\[
\sqrt{1-\Tr[UPU^*|c\rangle\langle c|]}
\le(2\sqrt{C_{\rm dec}}+\sqrt{50})\sqrt{e_B(P)}.
\]
Measure in the coordinate basis after $U$ and prepare $P_c$ on coordinate $c$; prepare an arbitrary representative on added coordinates. Its loss on good states is at most
\[
K_{\rm in}e_B(P),\quad
K_{\rm in}=(2\sqrt{C_{\rm dec}}+\sqrt{50})^2+50<2500.
\]
Outside the good set loss is at most $1<2800e_B(P)$. If that set is empty, any basis measurement with pure outputs suffices. This strengthens Theorem~\ref{thm:surrogate}: one surrogate can have \emph{rank-one} effects
\[
E_a=|u_a\rangle\langle u_a|,\quad u_a=U^*|a\rangle,
\qquad d\le N\le\lfloor8d/7\rfloor,\quad \sum_aE_a=I.
\]
The frame Gram matrix $UU^*$ is a rank-$d$ projection. In particular its collective off-diagonal leakage at each representative satisfies
\[
\sum_{b\ne c}|\langle u_b,u_c\rangle|^2
=(UU^*)_{cc}\bigl(1-(UU^*)_{cc}\bigr)
\le4C_{\rm dec}e_B(P_c).
\]
Thus the representation controls all other frame directions simultaneously. A coordinate-projective Naimark realization embeds $\mathbb C^N$ into $H\otimes\mathbb C^2$; one ancillary qubit suffices. This is a Hilbert-space statement, not a circuit-synthesis bound.

\subsection{Sharp order of the excess outcome budget}

Fix $0\le\varepsilon\le1/2800$, and construct a frame only for a family in $S_\varepsilon=\{P:e_B(P)\le\varepsilon\}$. Keep the classes that meet this family. Their representatives still have error at most $\varepsilon$, and
\[
k_\varepsilon(1-C_{\rm dec}\varepsilon)
\le\sum_c\Tr(M_cP_c)\le\sum_c\Tr M_c\le d.
\]
The selected compressed effects need only sum to at most $I$ for the polar proof. Hence a single frame works on the whole selected family, with loss $<2500e_B(P)$ and
\[
N_\varepsilon\le\min\left\{\lfloor8d/7\rfloor,
\left\lfloor\frac d{1-C_{\rm dec}\varepsilon}\right\rfloor\right\},\qquad
N_\varepsilon-d\le\frac{C_{\rm dec}d\varepsilon}{1-C_{\rm dec}\varepsilon}.
\tag{7a}
\]
If $C_{\rm dec}\varepsilon(d+1)<1$, then $N_\varepsilon=d$ and the measurement is one orthonormal basis. This restricted construction is not asserted to retain the all-pure guarantee outside the chosen family.

The order $d\varepsilon$ is necessary. Take $q+1$ real regular-simplex vectors in dimension $q$, with off-diagonal inner products $-1/q$. Put $\phi_i=\psi_i\otimes\psi_i$ and $t=(q+1)^{-1}\sum_i\phi_i$. Direct calculation gives
\[
\norm t^2=\langle t,\phi_i\rangle=1/q,\qquad
v_i=\frac{\phi_i-t}{\sqrt{1-1/q}},\qquad
\langle v_i,v_j\rangle=\langle\psi_i,\psi_j\rangle.
\]
Thus $V\psi_i=v_i$ extends to an isometry, and $B(P)=VPV^*$ has $e_B(P_i)=1/q$. Take $b$ orthogonal blocks: $d=bq$, $\varepsilon=1/q$, and there are $d+d\varepsilon$ exhibited inputs. Distinct inputs have overlap modulus at most $1/q$.

If a prepared density matrix $\sigma$ served two different inputs at fidelity $f>(1+1/q)/2$, it would violate
\[
2f<\Tr[\sigma(P_i+P_j)]\le\norm{P_i+P_j}_\infty\le1+1/q.
\]
Every input attaining fidelity $f$ needs some prepared state attaining at least $f$. Consequently any measure-and-prepare surrogate with loss at most $K\varepsilon$, for fixed $K$ and integer $q>\max(2800,2K+1)$, requires at least $d+d\varepsilon$ distinct preparations, and at least that many outcomes. This matches the order in (7a), allows mixed preparations, and shows that a single fixed basis measurement cannot always suffice. It does not exclude externally randomized choices among many bases.

For any prior, the same frame measurement also retains the pure ensemble's Holevo information up to $g_d(2800E)$, where $E=\sum_i p_i e_B(P_i)$ and $g_d(t)=h_2(t)+t\log_2(d-1)$ up to $t=1-1/d$, capped at $\log_2d$ thereafter. Indeed the posterior state on outcome $y$ has maximum eigenvalue at least its fidelity with the pure state prepared there. Its entropy is at most $g_d$ of that posterior loss. Concavity and data processing $I(Y:A)\le I(Y:X)$ give $0\le\chi-I(X:Y)\le g_d(2800E)$. The companion \texttt{proofs/frame\_reconstruction.txt} includes the full argument and numerical-tolerance contract.

The final mixed-state proof supplies a still stronger starting classifier for this pure-frame conversion: replace $C_{\rm dec}$ by $C_*<11$ and the class-count bound by $k<200d/199$. The identical polar proof then gives good-state loss $<250e_B(P)$, an all-pure rank-one frame with $N\le\lfloor200d/199\rfloor$, and $N_\varepsilon-d\le C_*d\varepsilon/(1-C_*\varepsilon)$ on a uniformly good family. The earlier constants and Schur proof are retained to expose the source mechanism and the route by which the stronger result emerged; they are not the final bottleneck.

\subsection{The amplification operation applies to mixed decision experiments}

The amplification step has a separate scope beyond pure-state reconstruction. Let $\{\rho_i\}$ be any finite mixed-state experiment, with supplied class labels $c(i)$, and assume one POVM has every wrong-class probability at most $a<1/2$. Put
\[
e_i=1-F(B(\rho_i),\rho_i\otimes\rho_i),\qquad
\beta(a)=3a-2a^2,\qquad
K(a)=\frac8{a(1-\sqrt{\beta(a)})^2},
\]
where $F$ is squared Uhlmann fidelity. There is one POVM, independent of the prior, whose wrong-class probability on every $i$ is at most $K(a)e_i$.

Indeed the Bures metric $d_B=\sqrt{2(1-\sqrt F)}$ is contractive, so the correlated two-call three-output tree is within $\sqrt{8e_i}$ of $\rho_i^{\otimes3}$. Uhlmann purifications imply $|\sqrt{\Tr Q\sigma}-\sqrt{\Tr Q\tau}|\le d_B(\sigma,\tau)$ for every effect $Q$. Ideal majority error is at most $f(r)=3r^2-2r^3\le\beta(a)r$ when $r\le a$. Set $\tau(a)=a(1-\sqrt{\beta(a)})^2/8$. For $e_i\le\tau(a)$ amplification remains feasible and constrained optimality gives average risk at most $8E/(1-\sqrt{\beta(a)})^2$. Charge the bad labels by $1/\tau(a)$ and apply minimax. This proves the assertion and controls every bounded class-dependent decision loss by the same error.

The weak-classifier premise is necessary for arbitrary supplied labels: identical states with different labels have an exact constant product replicator but cannot reveal a balanced binary label with error below $1/2$. Section~\ref{sec:mixed} closes the initially missing bridge by constructing the appropriate fidelity classes and their global weak measurement through deletion induction. Purifying the promise remains invalid: a replacer clones a single maximally mixed state exactly while destroying its canonical purification.

For a supplied pure ensemble let $G_{ij}=\sqrt{p_ip_j}C_{ij}$, and let $H$ be any positive semidefinite correlation matrix. If
\[
\norm{G-G\circ C\circ H}_1\le\varepsilon,
\]
realize $H$ as an environment Gram matrix. Uhlmann's theorem aligns the input and target column operators and gives a physical cloner with average global infidelity at most $\varepsilon$ (the sharper quantity is $1-F_{\rm root}(G,G\circ C\circ H)^2$). Hence
\[
e_{\EB}\le\min(1,2800\varepsilon),\qquad
\sum_i p_i\tfrac12\norm{\Lambda(P_i)-P_i}_1
\le\min(1,\sqrt{2800\varepsilon}).\tag{8}
\]
This answers the supplied finite-ensemble trace-Gram interface with a linear infidelity modulus. It does not assert mixed-state reconstruction, diamond-norm approximation, efficient acquisition, or quantum-capacity consequences.

\section{Sharp transport stability under one radial moment budget}

Let $\rho(x)\,dx$ be a full-dimensional log-concave probability with compact support $K\subset B(0,R)\subset\mathbb R^d$, $d\ge2$. Define finite slice constants
\[
A_j=\int_{e_j^\perp}\operatorname*{ess\,sup}_{t\in\mathbb R}\rho(z+te_j)\,dz,
\qquad A=\sum_jA_j.
\]
For targets $\mu,\nu\in\mathcal P_2$, let $T_\mu,T_\nu$ be the unique quadratic Brenier maps from this source. Put $w=W_2(\mu,\nu)$, $D=\norm{T_\mu-T_\nu}_{L^2(\rho)}$, and $\tau_\mu(L)=\int_{|y|>L}|y|^2d\mu(y)$.

\begin{theorem}[One-sided tail interpolation]\label{thm:transport}
For every $h,L>0$, with $c_P=1+\sqrt{162}$,
\[
D^2\le\frac{12d c_P^2R^2w^2}{h^2}
+128AL^2h+40d\tau_\mu(L)+80dw^2.\tag{9}
\]
In particular, if only $\int|y|^p d\mu\le M$, $p>2$, while $\nu$ is arbitrary in $\mathcal P_2$, then
\[
D\le w+Cw^{(p-2)/(3p-2)}.\tag{10}
\]
The small-distance exponent is sharp on the fixed uniform cube, even with two-atom targets satisfying the moment budget on both sides.
\end{theorem}

Here and below constants may depend on the displayed source, dimension, and budget; no dimension-independent constant or computational procedure is inferred. The full finite-cell and approximation proofs are in \texttt{proofs/transport.txt}. We give the actual load-bearing operations next.

\subsection{Weighted midpoint domination and potential control}

For convex sets $U,V$ of source masses $\alpha,\beta>0$, use the same sequential conditional-quantile parameter for their conditional source laws $X,Z$. Their triangular diagonal derivatives telescope to $\alpha/\rho(X)$ and $\beta/\rho(Z)$. Arithmetic--geometric mean and log-concavity imply for $H=(X+Z)/2$ that
\[
\det_{\rm triangular}DH
\ge\frac{\sqrt{\alpha\beta}}{\sqrt{\rho(X)\rho(Z)}}
\ge\frac{\sqrt{\alpha\beta}}{\rho(H)}.
\]
Successive one-dimensional substitutions therefore give the measure domination
\[
\sqrt{\alpha\beta}\operatorname{Law}((X+Z)/2)\le\rho(x)dx.\tag{11}
\]
If all midpoints lie in one central Laguerre cell, the domination is restricted to that cell.

Write $M_i,B_i$ for the source masses and first moments of Laguerre cells. Compare cells at opposite intercept perturbations and apply (11) to $R^2-|x|^2$. If their masses are $\alpha_i,\beta_i$, set $\gamma_i=\sqrt{\alpha_i\beta_i}$,
$\mathcal D=\sum_i\gamma_i\mathbb E|X_i-Z_i|^2$, and
$\mathcal B=\sum_i(\alpha_i-\beta_i)^2/\gamma_i$. The partitions and the quadratic midpoint identity give
\[
\mathcal D/4\le R^2\mathcal B/8+(R/2)\sqrt{\mathcal D\mathcal B},
\qquad\mathcal D\le5R^2\mathcal B.
\]
Passing to derivatives yields
\[
\sum_i\frac{|DB_i[q,0]|^2}{M_i}
\le162R^2\sum_i\frac{|DM_i[q,0]|^2}{M_i}.\tag{12}
\]
Interpolate the target sites along an optimal target coupling. Their cell masses are held fixed by intercepts. The mass derivative in intercepts is a connected weighted graph Laplacian; its adjoint cancellation with the moment derivative and (12) bounds the intercept speed by $\sqrt{162}Rw$. Consequently potentials $u,v$ can be normalized so that
\[
F:=\norm{u-v}_{L^2(\rho)}\le c_PRw.\tag{13}
\]
This argument bounds no inverse-Laplacian condition number. That distinction matters for any proposed numerical algorithm.

\subsection{Interpolation toward slice modes and one-sided tails}

On each coordinate slice shift a point by $h$ toward a mode of its unimodal source density, excluding a strip of width $2h$ about the mode. Density does not decrease under the shift. The two shifted images have disjoint interiors, so
\[
\int_{\rm allowed}|D_h(u-v)|^2\rho\le4F^2/h^2.
\]
At good points, both endpoint derivatives have magnitude at most a cutoff $L$. Their clipped monotone derivatives have total variation at most $2L$; integration against the slice peak bounds the squared finite-difference error by $4L^2hA_j$ for each potential. The excluded strip costs $8L^2hA_j$. At bad points the mass is at most $4\max(\tau_\mu(L),\tau_\nu(L))/L^2$, and splitting each squared gradient at radius $L$ costs at most five times that maximum tail. Thus
\[
D^2\le12dF^2/h^2+32AL^2h
+20d\max(\tau_\mu(L),\tau_\nu(L)).\tag{14}
\]
For an optimal coupling $(X,Y)$ of $\mu,\nu$, split $|Y|>L$ according to $|X|>L/2$. The inequalities $|Y|^2\le2|X|^2+2|Y-X|^2$ on the first part and $|Y|\le2|Y-X|$ on the second give
\[
\tau_\nu(L)\le2\tau_\mu(L/2)+4w^2.\tag{15}
\]
Equations (13)--(15), using cutoff $2L$, prove (9). For a $p$th moment choose
$L=w^{-2/(3p-2)}$, $h=w^{2p/(3p-2)}$ when $w\le1$. All relevant powers are $w^{2(p-2)/(3p-2)}$. For $w>1$, $D\le w+2M^{1/p}$. Finite norm-decreasing target quantizers and smooth compact log-concave source convolutions remove the preliminary smooth/finite restrictions; optimal-plan uniqueness and quadratic uniform integrability justify the limits.

\subsection{An exact sharp modulus for general radial budgets}

Let $\Phi:[0,\infty)\to[0,\infty)$ be finite and nondecreasing, $\Phi(0)=0$, with $\Phi(L)/L^2$ nondecreasing and tending to infinity. Convexity of $\Phi$ is unnecessary. For $M>0$, the one-sided class modulus on a fixed cube source is
\[
\Omega_{\Phi,M}(w)=\sup\{\norm{T_\mu-T_\nu}_2:
\int\Phi(|y|)d\mu\le M,\ \nu\in\mathcal P_2,\ W_2(\mu,\nu)\le w\}.
\]
Choose fixed $L_0$ with $M/\Phi(L_0)<1/4$ and define
\[
Q(w)=\inf_{L\ge L_0}\left\{w^{2/3}L^{4/3}+\frac{ML^2}{\Phi(L)}\right\}.
\]
The complete sharp statement, for sufficiently small $w$, is
\[
\frac18Q(w)\le\Omega_{\Phi,M}(w)^2\le C Q(w).\tag{16}
\]
The upper bound follows by $\tau_\mu(L)\le ML^2/\Phi(L)$ in (9) and optimizing $h$; $w^2$ is absorbed since $Q(w)\ge w^{2/3}L_0^{4/3}$.

For the lower bound, on the uniform cube put $c=1-2a$, $0<a<1/4$, and use potentials
\[
u(x)=\max(0,L(x_1-c)),\quad
v(x)=\max\left(0,\frac{L(x_1-c+ax_2)}{\sqrt{1+a^2}}\right).
\]
Their outlier cells have mass $a$, symmetric-difference mass $a/4$, and intersection mass $7a/8$. Both nonzero target atoms have radius exactly $L$. Exact calculations give
\begin{align*}
w^2&=2aL^2(1-(1+a^2)^{-1/2})\le a^3L^2,\\
D^2&=aL^2/4+(7a/8)\,2L^2(1-(1+a^2)^{-1/2})\ge aL^2/4.
\end{align*}
Set $a=\min\{M/\Phi(L),(w/L)^{2/3}\}$. Both moment budgets hold and the target distance is at most $w$. Thus $\Omega^2\ge\frac14\sup_L\min(A(L),B(L))$, where $A=w^{2/3}L^{4/3}$ is continuous increasing and $B=ML^2/\Phi(L)$ is nonincreasing. If this supremum is $S$, choose $A(L)=S+\epsilon$; necessarily $B(L)\le S$. Hence $Q\le2S$, proving (16), even for jump discontinuities in $\Phi$.

For $\Phi(L)=L^p[\log(e+L)]^r$, $p\ge2$, $r\ge0$, with $p>2$ or $r>0$, the sharp order is
\[
\Omega(w)\asymp
w^{(p-2)/(3p-2)}[\log(1/w)]^{-2r/(3p-2)}.
\tag{17}
\]
At $p=2,r>0$ this is $[\log(1/w)]^{-r/2}$. A bound on second moments alone has no vanishing class-uniform modulus: take $L=a^{-1/2}$ in the two-atom family, so $w\to0$ while $D^2\ge1/4$. More generally, under the ratio-monotonicity hypothesis, the radial-budget class is uniformly continuous exactly when $\Phi(L)/L^2\to\infty$ (the degenerate zero-budget function also fails). This is class-uniform continuity, not a denial of qualitative stability at a fixed target limit.

\subsection{Scope against the closest checked primary results}

\href{https://arxiv.org/html/2610.09640v1}{Wang and Wang, Quantitative stability of optimal transport via regularization}, Theorem 2, permits broader BV-reference and unbounded sources. On the common bounded-density source interface, its one-sided $p$th-moment exponent is $(p-2)/(4p-2)$ against $W_2$. Equation (10) has the stronger exponent $(p-2)/(3p-2)$ on compact log-concave sources. It does not subsume their broader source class.

\href{https://arxiv.org/html/2610.07577v1}{Ruan, Sharp conditioning and quantitative stability of optimal transport}, Theorem 10 and Appendix A.1, give a related smaller map exponent against $W_1$; its power-log conditioning theorem uses an objective gap, a different input. The second-moment failure is already present there and is not claimed new in isolation. Other exact comparisons, including Delalande--M\'erigot, Letrouit--M\'erigot, Mischler--Trevisan, and M\'erigot's May 2026 manuscript, are recorded in the companion proof. No inspected statement matches (16); this is a limited comparison, not a priority certificate.

The compact-source assumption is load-bearing in this proof. For a fixed Gaussian source and cells separated at coordinate $t$, $DM_2=\varphi(t)$ and $DB_2=t\varphi(t)e_1$, with opposite derivatives in cell one. The ratio of the two quadratic forms in (12) equals $t^2$. It diverges while every source moment stays fixed. Thus a plain replacement of $R^2$ by a source second moment is false.

\subsection{A centered operation removes bounded support under curvature}

There is nevertheless a complete stronger extension, using different mathematics. Suppose
\[
\rho(x)=e^{-\kappa|x|^2/2}h(x),\quad \kappa>0,
\]
where $h$ is nonnegative log-concave. The convex support may be bounded, a proper unbounded set, or all of $\mathbb R^d$. The slice constants $A_j$ are finite by a Gaussian envelope. Then (9) holds with its first coefficient replaced by $24d/\kappa$; the one-sided sharp moment and radial-budget upper bounds all follow.

Here is the essential new centered argument. Strong log-concavity improves the midpoint substitution to
\[
M_i\ge\gamma_i\mathbb E\exp(\kappa|X_i-Z_i|^2/8).
\]
Since $\sum_i(M_i-\gamma_i)\le\tfrac18\sum_i(\alpha_i-\beta_i)^2/\gamma_i$, it follows that
\[
\sum_i\gamma_i\mathbb E|X_i-Z_i|^2
\le\kappa^{-1}\sum_i\frac{(\alpha_i-\beta_i)^2}{\gamma_i}.
\]
With $\bar x_i=B_i/M_i$, passage to derivatives yields
\[
\sum_i\frac{|DB_i-\bar x_iDM_i|^2}{M_i}
\le\kappa^{-1}\sum_i\frac{DM_i^2}{M_i}.\tag{21}
\]
For moving target sites set $a_i=r_i+\bar x_i\cdot v_i-\int\partial_tu_t\rho$. Choose $DM_i[q]=M_i a_i$. The centered adjoint identity is
\[
\sum_i M_i a_i^2=-\sum_i v_i\cdot(DB_i[q]-\bar x_iDM_i[q]).
\]
Equation (21) bounds this by $w\norm a_M/\sqrt\kappa$. Conditioning a strongly log-concave density on a convex cell preserves the classical covariance bound $\operatorname{Cov}_{P_i}(x)\preceq\kappa^{-1}I$. Therefore
\[
\norm{\partial_tu_t-\int\partial_tu_t\rho}_2^2
=\sum_iM_i\operatorname{Var}_{P_i}(x\cdot v_i)+\sum_iM_i a_i^2
\le2w^2/\kappa.
\]
Source-mean-normalized potentials satisfy $F\le\sqrt{2/\kappa}\,w$. The slice interpolation and one-sided tail transfer now give
\[
D^2\le\frac{24d}{\kappa}\frac{w^2}{h^2}+128AL^2h+40d\tau_\mu(L)+80dw^2.\tag{22}
\]
The covariance premise follows directly by completing the square in the log-Laplace transform: $\log Z(b)-|b|^2/(2\kappa)$ is concave by Gaussian convolution log-concavity. Compact truncations preserve $\kappa$ and have $A_j(\rho_R)\le A_j(\rho)/\rho(B_R)$; hence the radius-free estimate passes to the unbounded source. Gaussian smoothing has curvature $\kappa/(1+\kappa\epsilon)$ and does not increase $A_j$, which handles nonsmooth sources. The complete limit argument is in \texttt{proofs/strong\_transport.txt}.

Sharpness holds even for one fixed genuinely unbounded source, the product of a Gaussian restricted to $[-1,1]$ in the first coordinate and unrestricted Gaussians in the others. Tilt a rare boundary cap by angle $\arctan a$, adjusting its intercept to keep mass $a$. Its symmetric difference with the original cap is asymptotic to a positive constant times $a$, whereas equal-radius target atoms have $w\asymp La^{3/2}$ and map distance $D\asymp L\sqrt a$. Dominated convergence includes the rare slices where the affine threshold exits $[-1,1]$. This proves the same sharp radial-budget profile on that fixed unbounded slab. It does not prove sharpness for a source with full support $\mathbb R^d$.

\section{A second complete consequence from orbit geometry}

The companion \texttt{proofs/orbit\_memory.txt} proves a general transfer. If an orthogonal compact-group action on $V$ satisfies, uniformly for all nonzero $x$ and all centers $y$,
\[
\Pr_g(\norm{gx-y}\le b)\le\min\{1,C(b/\norm{x})^k\},
\]
and a fixed finite gate distribution has an absolute regular $L^2_0$ spectral gap, then arbitrary time-dependent gate-aware positive $m$-state lifts satisfy
\[
F_m(T)\le2\exp[-cT/m^{2/k}].\tag{18}
\]
The mechanism is matched-Gaussian relative entropy: for a conditional-mean law $\mu$, Gaussian smoothing against the same Gaussian reference gives a potential decreasing under convex-order merging, while the gate gap spends potential on its noninvariant part. Finite support and the cap bound force such a part whenever retained second moment is positive. A half-life argument and submultiplicativity give (18), without a Haar-approximation logarithm.

For conjugation on traceless Hermitian $d\times d$ matrices, a degeneracy-uniform cap bound has $k=2d-2$. Some adjacent eigenvalue gap is at least $\norm H_{\rm HS}/[\sqrt d(d-1)]$; its spectral projection reduces the cap to a complex Grassmann cap of real dimension $2j(d-j)\ge2d-2$. Rank-one orbits show this exponent is optimal. A pure-projector net supplies the matching positive interpolation upper bound.

For an explicit fixed alphabet of algebraic adjacent SU(2) gates in SU($d$), and the target terminal probability
\[
\Tr\left[E G_t\left(\rho/2+I/(2d)\right)G_t^*\right],\qquad t\le T,
\]
the exact supplied interface has optimal classical state cost
\[
m_*(T;\epsilon)=\Theta_d((T+1)^{d-1}),\quad
0\le\epsilon\le[8d(d^2-1)]^{-1}.\tag{19}
\]
Only an $m$-valued register is persistent; time and deadline are supplied, input preparation is provided at the start, and a terminal binary effect is provided at the end. No stored word, preparation label, external seed, or accumulated matrix is free. The lower bound uses only $2(d^2-1)$ preparations and one fixed binary query. The upper is exact but its read-only transition descriptions and compiler costs are separately charged. The group gap is the imported \href{https://arxiv.org/abs/1108.6264}{Bourgain--Gamburd SU($d$) theorem}; no numerical all-frequency gap is certified.

Allowing varying state counts does not allow occasional large-memory steps to recharge the potential. For any fixed positive terminal retained correlation,
\[
\#\{t<T:m_t\le M\}\le K M^{1/(d-1)},\qquad
\sum_{t<T}\log_2m_t\ge(d-1)\log_2(T!)-O_d(T).\tag{20}
\]
The source's qubit predecessor is credited to Astra N565--N570. The higher-dimensional cap transfer and (20) are the present deductions. They are fixed-$d$ memory laws, not a total-device computational advantage.

\subsection{Minimum orbit dimension classifies the positive composition gap}

The geometry generalizes beyond Hermitian conjugation. Let a compact Lie group act orthogonally on a finite-dimensional real $V$, and suppose
\[
k=\min_{x\ne0}\dim(Gx)>0.
\]
Under the same full regular group gap, the joint spectral radius of all positive $m$-state lifts obeys
\[
1-\rho_m=\Theta(m^{-2/k}).\tag{23}
\]
Here $\rho_m=\lim_{T\to\infty}F_m(T)^{1/T}$. Constants depend on the action and the gate law. The matching lower lift may live in the invariant span of a minimum-dimensional orbit; (23) alone does not supply an upper simulator for arbitrary inputs outside that span.

For completeness, the two geometric interfaces can be derived directly. The differential $L_x:\mathfrak g\to V$, $L_x\xi=d\rho(\xi)x$, has rank at least $k$ for every unit $x$. Its $k$th singular value has a positive minimum on the compact unit sphere. Uniform exponential charts, projection onto $k$ nonsingular image directions, and change of variables give a cap bound $C(b/\norm x)^k$ uniformly in $x$ and the cap center. The matched-Gaussian argument then gives the upper bound (18).

For the match choose a minimum orbit $M$ of dimension $k$, $W=\operatorname{span}M$, and $K=\operatorname{conv}M$. There are no fixed vectors, so the Haar mean is zero and $rB_W\subset K$ for some $r>0$. A geodesic $\delta$-net has $O(\delta^{-k})$ vertices. At a maximizer of any support functional its first derivative vanishes; the bounded second fundamental form therefore makes its support loss $O(\norm H\delta^2)$. Thus $(1-C\delta^2)K\subset\operatorname{conv}\{v_i\}$. Stochastic barycentric interpolation gives a dual positive-lift eigenvalue $1-C\delta^2$. Taking $\delta\asymp m^{-1/k}$ proves the other side of (23). The detailed uniform-chart and eigenvalue arguments are in \texttt{proofs/minimum\_orbit.txt}.

\subsection{One fixed two-qubit normalized trace experiment}

For unit quaternions $q\in S^3$, write $U(q)=q_0I-i\sum_{j=1}^3q_j\sigma_j$. Left multiplication is orthogonal on $\mathbb R^4$ and has orbit dimension three. Define
\[
\varrho_q=\frac14\begin{pmatrix}I_2&U(q)^*/2\\U(q)/2&I_2\end{pmatrix},\quad
E_0=\frac12\begin{pmatrix}I_2&I_2\\I_2&I_2\end{pmatrix},\quad
C(g)=\operatorname{diag}(I_2,U(g)).
\]
The preparations have eigenvalues $3/8,1/8$, each twice; $E_0$ is a projector. Starting from the single fixed state $\varrho_{e_0}$ and applying a word $g_1,\ldots,g_t$ gives
\[
p_*(+\mid g_1,\ldots,g_t)
=\frac12+\frac18\operatorname{Re}\Tr[U(g_t)\cdots U(g_1)].\tag{24}
\]
Use the fixed algebraic SU(2) gap alphabet and also the eight quaternion gates $\{\pm I,\pm i\sigma_x,\pm i\sigma_y,\pm i\sigma_z\}$. For all words of length at most $T$ and fixed absolute probability error $0\le\epsilon\le1/32$, the optimal persistent classical state count is
\[
m_*(T;\epsilon)=\Theta((T+1)^{3/2}).\tag{25}
\]
To derive the lower bound, use the first gate to realize the eight signed quaternion-coordinate preparations. Thereafter independent gap-alphabet gates preserve their isotropic second moment $I_4/4$. A fixed-query terminal score is $1/8$ and remains at least $1/16$ under the allowed error. Apply (18) with $k=3$. The prefix reduction charges one step and permits one fixed preparation and one fixed effect throughout.

For the upper bound an explicit rational stereographic net on $S^3$ has $N=2(2n+1)^3$, $n=\lceil\sqrt{3(T+1)}\rceil$, and support loss $3/(2n^2)$. Inflating the net and compensating its stochastic damping at the terminal readout is exact up to horizon $T$; the readout stays in $[0,1]$ by Bernoulli's inequality. The finite compiler costs $O(|\mathcal A|N^6\operatorname{poly}(\log(T+2)))$ bit operations by enumeration of five-vertex bases, with sparse algebraic rows. Terminal exact arithmetic can use $O(T\log(T+2))$ temporary bits. These descriptions and temporary costs are not part of the persistent gate-history register and are explicitly charged; (25) is not a claim of a total-memory or time speedup. The active coherence representation distinguishes $U$ and $-U$, which the Bloch adjoint representation identifies. That exact change of observable explains the exponent $3/2$ in place of the bare-qubit exponent one.

The same fixed probability process has a full-rank qutrit realization. Put $\psi_q=U(q)e_1\in\mathbb C^2$ and
\[
A_q=\begin{pmatrix}0&\psi_q^*\\\psi_q&0_2\end{pmatrix},\quad
R_q=I_3/3+A_q/4,\quad F_u=I_3/2+A_u/2.
\]
Their spectra are respectively $\{7/12,1/3,1/12\}$ and $\{1,1/2,0\}$; gates $\operatorname{diag}(1,U(g))$ act by quaternion multiplication. Since $\Tr(A_uA_q)=2\langle u,q\rangle$, $\Tr(F_uR_q)=1/2+\langle u,q\rangle/4$, exactly the same process. Qutrit dimension three is minimal at error $\epsilon\le1/32$: the eight signed-coordinate prefixes and conjugate suffixes give the matrix $M=\tfrac12\mathbf1\mathbf1^T+\tfrac14BB^T$, with $B$ having rows $\{\pm e_j\}_{j=0}^3$. Its five nonzero eigenvalues are $4,1/2,1/2,1/2,1/2$. Any qubit-memory probability matrix factors through the four-dimensional real space of Hermitian qubit operators, hence has rank at most four. Its operator-norm error is at least $1/2$, while entrywise error $\epsilon$ bounds that norm by $8\epsilon$. Thus a qubit requires $\epsilon\ge1/16$ already for these length-two words. This uses the same counted-memory contract and grants no uncounted side memory.

\section{Full-support sources: a sharp one-third law with one higher moment}

The previous radial profile is sharp over its class, but boundary geometry conceals a stronger operation on smooth full-support sources. Let $\gamma_d$ be standard Gaussian, let $p>2$, and assume only $\int|y|^p d\mu\le M$, while $\nu\in\mathcal P_2$ is arbitrary. Then
\begin{theorem}[Gaussian one-sided higher moment]\label{thm:gaussian}
For $d\ge2$,
\[
\norm{T_\mu-T_\nu}_{L^2(\gamma_d)}\le w+C_{d,p,M}w^{1/3},\qquad w=W_2(\mu,\nu).
\tag{G1}
\]
The exponent $1/3$ is sharp on this single fixed full-support source, even with both targets supported in one compact ball. The constant is not uniform as $p\downarrow2$.
\end{theorem}

Here is the additional mechanism, with its exact interface. For $a\le b$,
\[
(b-a)^2\le2(b|b|-a|a|).
\tag{G2}
\]
Let $u,v$ be finite-target potentials normalized by their Gaussian means. The centered potential estimate proved above gives $F=\norm{u-v}_2\le\sqrt2w$. Clip the monotone coordinate derivatives at $[-L,L]$, writing $g_{z,j}=\operatorname{clip}(\partial_jz,[-L,L])$ for $z=u,v$. Shift a coordinate $t$ to $t+h$ when $t\le0$ and to $t-h$ otherwise, with $0<h\le1$. The Gaussian density satisfies $r(t)\le cr(S(t))$, $c=e^{1/2}<5/3$; the shift has multiplicity at most two. Consequently its difference quotient satisfies $\int|D_h(u-v)|^2d\gamma_d<10F^2/h^2$.

At points where the endpoint derivatives stay below $L$, convexity bounds $|D_hz-\partial_jz|$ by $|g_{z,j}(S(t))-g_{z,j}(t)|$. Apply (G2) and the weighted Stieltjes inequality
\[
\int|H(S(t))-H(t)|r(t)dt\le2ch\int r(t)dH(t)
\]
to $H=g_{z,j}|g_{z,j}|$. The inequality follows by Tonelli on intervals of length $h$; the density distortion is at most $c$. Boundedness of $H$ removes the boundary term in Gaussian integration by parts. With $q=p/(p-2)$ and $m_q=(\mathbb E|G|^q)^{1/q}$,
\[
\int_{\rm good}|D_hz-\partial_jz|^2d\gamma_d
\le4ch\,\mathbb E[X_jg_{z,j}|g_{z,j}|]
\le4ch\,m_q\norm{g_{z,j}}_p^2.
\tag{G3}
\]
The bad set has mass below $9\tau(L)/L^2$, with $\tau(L)=\max(\tau_\mu(L),\tau_\nu(L))$, and derivative-difference energy at most $40\tau(L)$ per coordinate. Summing and using the three-term squared triangle gives
\[
D^2\le30dF^2/h^2+
20h m_q\sum_j(\norm{g_{u,j}}_p^2+\norm{g_{v,j}}_p^2)
+40d\tau(L).
\tag{G4}
\]
This replaces the former $L^2h$ cost by an $h$ cost controlled by clipped moments.

For an optimal target coupling $(Y,Z)$, a pointwise split at $|Z|=2|Y|$ gives
\[
\mathbb E\min(|Z|,L)^p\le2^pM+4L^{p-2}w^2.
\tag{G5}
\]
It grants no higher moment to $\nu$. For $0<w\le1$, choose $L=w^{-2/(p-2)}$. Then $\norm{g_{u,j}}_p^p\le M$, $\norm{g_{v,j}}_p^p\le2^pM+4$, and the one-sided tail transfer gives $\tau(L)\le(2^{p-1}M+4)w^2$. In (G4) put $h=w^{2/3}$. This proves (G1), for example with
\[
C_{d,p,M}=\max\left\{2M^{1/p},
\sqrt{60d+20dm_q[M^{2/p}+(2^pM+4)^{2/p}]+40d(2^{p-1}M+4)}\right\}.
\]
For $w>1$, use $D\le w+2M^{1/p}$. Finite norm-decreasing target quantization preserves the one-sided budget. Common-source optimal-plan uniqueness and quadratic uniform integrability pass the result to arbitrary admitted targets. The full argument, including finite-potential limits, is in \texttt{proofs/gaussian\_transport.txt}.

\paragraph{Exact sharpness.}
Let $X,Y$ be independent scalar standard Gaussians, $H(s)=\mathbb P(|X|<s)$ for $s\ge0$, and $H(s)=0$ otherwise. Choose $C_a$ by
\[
\mathbb E H(a(C_a+Y/2))=H(a)=:m_a.
\]
Monotonicity gives a unique $C_a$, and dominated convergence gives $C_a\to C_*$ where $\mathbb E(C_*+Y/2)_+=1$. The potentials
\[
u_a=\max\{-x_1,x_1,a\},\qquad
v_a=\max\{-x_1,x_1,aC_a+(a/2)x_2\}
\]
have exactly equal central-cell masses and, by $x_1$ symmetry, exactly equal outer masses. Their gradient laws differ only by moving their central atom from $0$ to $(a/2)e_2$. Thus $w_a^2=m_a a^2/4\asymp a^3$ exactly. Their central cells have symmetric difference
\[
\mathbb E|H(a)-H(a(C_a+Y/2))|\asymp a,
\]
because after division by $a$ its limit is
$2\varphi(0)\mathbb E|1-(C_*+Y/2)_+|>0$.
The two coordinate components of the map difference are orthogonal, giving $D_a^2=\gamma_d(A_a\triangle B_a)+m_a a^2/4\asymp a$. Hence $D_a\asymp w_a^{1/3}$. All slopes stay in one fixed ball. A fixed common rescaling fits any prescribed $M>0$.

\paragraph{Score-moment extension without an upper curvature bound.}
Let $\rho=e^{-V}$ on all of $\mathbb R^d$, normalized to mass one, with $V\in C^1$, $V(x)-\kappa|x|^2/2$ convex, $\kappa>0$, and
\[
S_q=\sum_j\norm{\partial_jV}_{L^q(\rho)}<\infty,\qquad q=p/(p-2).
\]
The same one-sided target promise gives exponent $1/3$. Here is why no upper Hessian bound is necessary. On a full-line unimodal slice density $r$ tending to zero at both ends, any bounded monotone $g$ satisfies
\[
r(t)g(t)^2\le J:=\int|r'(s)|g(s)^2ds.
\tag{G6}
\]
If $g(t)\ge0$, integrate on the right tail starting at $\max(t,\text{mode})$; if $g(t)<0$, use the left tail ending at $\min(t,\text{mode})$. Monotonicity of $g$ and of the density on that tail proves (G6).

Use only shifts of length $h$ toward the mode that do not cross it, and exclude the strip of width $2h$. Density then never decreases, and shifted images do not overlap. The finite-difference term costs $12dF^2/h^2$. Quadratic variation, by (G2) and Stieltjes integration by parts, costs at most $6h\sum_j(J_{u,j}+J_{v,j})$. Equation (G6) bounds the strip's good-derivative energy by $4h\sum_j(J_{u,j}+J_{v,j})$. One unified bad set has mass at most $4\tau(L)/L^2$ per coordinate and energy at most $20\tau(L)$; it includes high derivatives inside the strip. Hence
\[
D^2\le12dF^2/h^2+10h\sum_j(J_{u,j}+J_{v,j})+20d\tau(L).
\tag{G7}
\]
H\"older gives $J_{z,j}\le\norm{\partial_jV}_{L^q(\rho)}\norm{g_{z,j}}_p^2$. The centered bound supplies $F^2\le2w^2/\kappa$. With the same cutoff $L=w^{-2/(p-2)}$ and $h=w^{2/3}$ this gives (G1), replacing the Gaussian constant by the square root of
\[
24d/\kappa+10S_q[M^{2/p}+(2^pM+4)^{2/p}]+20d(2^{p-1}M+4),
\]
enlarged to at least $2M^{1/p}$. Gaussian sources show the exponent is optimal over this class. The complete independent derivation is in \texttt{proofs/score\_transport.txt}. The source score and strong curvature are exact supplied assumptions; neither is inferred from the target moment. Proper-support densities can contribute a singular boundary variation not controlled by this score norm, consistent with their different sharp radial profile.

\subsection{A sharp logarithmic-moment phase on one full Gaussian source}

The score operation also resolves the critical moment scale. Fix $d\ge2$, $r>0$, $M>0$, and put $\Phi_r(t)=t^2\log(e+t)^r$. If just one target satisfies $\int\Phi_r(|y|)d\mu\le M$, while $\nu\in\mathcal P_2$ is arbitrary, then
\begin{theorem}[Gaussian logarithmic phase]\label{thm:logphase}
For the standard Gaussian source,
\[
\norm{T_\mu-T_\nu}_{L^2(\gamma_d)}
\le w+C_{d,r,M}w^{\beta_r},\qquad
\beta_r=\min\{1/3,r/(r+1)\},\quad w=W_2(\mu,\nu).
\tag{G8}
\]
This exponent is sharp for the entire small-distance class modulus, even when both targets satisfy the budget. At $r=0$ a uniform second-moment budget gives no vanishing class modulus on this same source.
\end{theorem}

Here is the exact balancing mechanism. If $Z\le L^2$ is any nonnegative source function, with no independence assumption from $X_j$, put $B_L=\log(e+L)$ and $s=(1/2-r)_+$. Splitting at $|X_j|=3\sqrt{\log(e+\sqrt Z)}$ gives
\[
\mathbb E|X_j|Z\le5\{1+\mathbb E Z\log(e+\sqrt Z)^r\}B_L^s.
\tag{G9}
\]
On the complementary event $Z\le e^{2X_j^2/9}$, and $\mathbb E|G|e^{2G^2/9}=18/(5\sqrt{2\pi})<2$. On the first event the asserted logarithmic factor controls the square root of the logarithm. This proves (G9) directly.

An optimal target coupling gives, with $a_r=4(1+\log2)^r$,
\[
\mathbb E\Phi_r(\min\{|Y_\nu|,L\})\le a_rM+4B_L^r w^2.
\]
Use the two cases $|Y_\nu|\le2|Y_\mu|$ and its complement. Also $\tau_\nu(L)\le2\tau_\mu(L/2)+4w^2$. With $L=e^B$, $B\ge1$, and $b_r=2(1-\log2)^{-r}$, (G7)--(G9) therefore give
\[
\begin{split}
D^2\le{}&24d w^2/h^2
+50d\,2^s hB^s\{2+(1+a_r)M+4\cdot2^rB^r w^2\}\\
&+20db_rM B^{-r}+80dw^2.
\end{split}
\tag{G10}
\]
For $0<r<1/2$ take $B=w^{-2/(r+1)}$, $h=w^{1/(r+1)}$; for $r\ge1/2$ take $B=w^{-2/(3r)}$, $h=w^{2/3}$. These give respectively $D^2=O(w^{2r/(r+1)})$ and $O(w^{2/3})$ for $w\le1$. For $w>1$, $D\le w+2\sqrt M$. Norm-decreasing finite quantizers preserve the supplied moment, and common-source graph-plan uniqueness plus uniform integrability passes the inequality to general targets. The exponential cutoff is a proof parameter; no bounded-bit algorithm is asserted.

For sharpness let $a_t=Q(t)=\mathbb P(G>t)$, $\theta_t=1/t$, and
\[
L_t=c a_t^{-1/2}\{\log(1/a_t)\}^{-r/2}.
\]
The two convex hinge potentials with slopes $q_t=L_te_1$ and $z_t=L_t(\cos\theta_t e_1+\sin\theta_t e_2)$ select equal-mass caps $A_t=\{X>t\}$ and $B_t=\{X\cos\theta_t+Y\sin\theta_t>t\}$. The targets put mass $a_t$ at these slopes and the remaining mass at zero. Their common budget tends to $c^2 2^{-r}$, so a fixed small $c$ fits every $M>0$. Gaussian tail change of variables and dominated convergence yield
\[
\frac{\mathbb P(A_t\cap B_t)}{a_t}
\longrightarrow\mathbb E e^{-(1/2-Y)_+}=2Q(1/2),\qquad
\frac{\mathbb P(A_t\triangle B_t)}{a_t}\longrightarrow2-4Q(1/2)>0.
\]
Indeed the conditional threshold is $t+s_t/t$ with $s_t\to1/2-Y$, and
$Q(t+s/t)=\varphi(t)t^{-1}\int_s^\infty e^{-z-z^2/(2t^2)}dz$ for $s\ge0$; the normalized intersection integrand is bounded by one. The optimal two-atom matching and the exact common-source map formula give
\[
w_t^2=2a_tL_t^2(1-\cos(1/t)),\qquad
D_t^2=L_t^2\mathbb P(A_t\triangle B_t)+|q_t-z_t|^2\mathbb P(A_t\cap B_t).
\]
Thus $D_t\asymp t^{-r}$, $w_t\asymp t^{-(r+1)}$. The distance is continuous and strictly decreasing at large $t$, so this proves a full small-distance lower modulus, rather than only a subsequence, for $r\le1/2$. For $r\ge1/2$ the compact three-atom construction above supplies the matching $1/3$ lower bound. At $r=0$, $D_t$ stays positive while $w_t\to0$. Complete endpoint and limit arguments are in \texttt{proofs/gaussian\_log\_phase.txt}, with a separate independent reconstruction in \texttt{proofs/gaussian\_log\_independent.txt}.

\paragraph{Primary comparison and its limit.}
The directly examined \href{https://arxiv.org/html/2610.09640v1}{Wang--Wang, 7 October 2026}, Theorem 2, covers a broader weighted-BV source class and gives the one-sided $p$-moment exponent $(p-2)/(4p-2)$ for Gaussian sources. It does not state (G1). \href{https://arxiv.org/html/2411.04908v3}{Letrouit--M\'erigot}, Theorem 1.4, assumes compact targets and gives weaker map exponents under strong curvature assumptions. These comparisons establish a difference from the checked statements, not exhaustive historical originality. The lower thin-strip mechanism is adapted from OpenAI family 374 and is credited accordingly.

\section{Other mathematics retained and exact failed interfaces}

\paragraph{Finite commutative replica products.}
For a finite commutative semigroup $S$, arbitrary input support $A$, and $m=|S|$, local eventual idempotents $e_a$ define monoids $I_a=e_aS$, units $u_a=e_aa$, and normalized defect monoids $J_a=\langle u_a^{-1}e_ab:b\in A\rangle$. A shortest defect word has at most $m-1$ letters. For $n\ge\max\{2m-1,\lceil(m-1)/\alpha\rceil\}$, $\alpha=\min_{a\in A}\mu(a)$, exact-equal-marginal product laws are within TV Hausdorff distance $(m-1)/n$ of transport supported on $z\in u_a^nJ_a$. This equals the idempotent relation $e_az=z$ together with $z\in A^n$. The correction preserves the input marginal exactly. Fixed exact target membership is eventually periodic in $n$, including zero supports; the full vertex and Cramer-sign proof is in the companion file. The prior source N674 already has the commutative binary-horizon LP, which is not claimed as a new result here.

\paragraph{The exact noncommutative boundary.}
For a generated finite completely simple semigroup in normalized Rees form $T=I\times G\times\Lambda$, word length induces a cyclic phase homomorphism $\phi:T\to\mathbb Z/h$; sandwich entries have phase zero. A finite block length $\ell\le2|T|^2$ fills the neutral phase. For $n\ge2+\ell\lceil1/\alpha\rceil$, the product-law polytope is exactly the laws $\nu$ with row marginal $\mu_I$, column marginal $\mu_\Lambda$, and support in $\phi^{-1}(n)$. Conditional endpoint coupling and an ordered repair block prove sufficiency. The companion \texttt{proofs/noncommutative\_repair.txt} gives the construction and cutoff proof.

General finite semigroups behave differently. There exist a finite two-generator semigroup, fair marginal $\mu=(\delta_{a_0}+\delta_{a_1})/2$, and a fixed nonzero annihilating target $z$ whose complete feasible-horizon set is
\[
\{70+60t:t\ge0,\ \operatorname{Re}[((-1+3i\sqrt3)/8)^t]\ge0\}.
\tag{S1}
\]
It is not eventually periodic. The construction uses fair checker bits $C_j=B\mathbin\oplus\mathbf1_{J=j}$. Their required fairness forces $B$ to be conditionally fair given the current finite state $J$, so its distribution evolves by a fixed Markov matrix even though the coupling may have arbitrary history dependence. An explicit 14-state chain realizes (S1)'s damped irrational rotation. Manchester encoding with balanced triple-run markers turns the accepted regular language into an infix code; all accepted words then form one syntactic-semigroup element $z$. This creates a single fixed product target rather than an accepted subset. The known nonperiodic Markov mechanism is credited in \texttt{proofs/nonperiodic\_semigroup.txt}; the checker/single-target coupling transfer is the present deduction. Infeasible horizons still admit output TV error at most $(\sqrt7/4)^t/2$, so this is an exact-attainment obstruction.

The same transfer works for any rational stochastic matrix $P$ and rational threshold $0<c<1$: a uniform alphabet with a common denominator and checkers $C_j=B+\mathbf1_{J=j}$ force conditional uniformity. Two terminal uniform-symbol tests force both $p_t(F)\ge c$ and $p_t(F)\le c$. Balanced wrapping therefore gives one finite semigroup with fixed marginal and point-mass target whose full feasible set is
\[
\{2(d+1)t+12:t\ge0,\ (\delta_sP^t)(F)=c\}.
\tag{S2}
\]
By the imported \href{https://people.mpi-sws.org/~joel/publications/markov_reach14.pdf}{Akshay--Antonopoulos--Ouaknine--Worrell, Theorem 3}, deciding arbitrary rational Markov exact reachability would decide the Skolem problem. Composing with (S2), any algorithm deciding whether \emph{some} feasible horizon exists for arbitrary finite-semigroup tables would decide that problem too, already with a uniform marginal and one target atom. The source reduction uses threshold $1/4$, so allowing time zero introduces no extra solution. The semigroup table can grow enormously: this is a computable reduction, with no polynomial table-size or undecidability claim. Full constructions and the imported-premise boundary are in the companion proof.

\paragraph{Cyclic noncommutative coupling.}
A labelled explicit finite graph admits a rational LP of size $O(N^3\log(n+1)+E)$ for laws on length-$n$ closed-walk words with every coordinate having a prescribed common label marginal. Balanced product trees carry expected production counts; random cyclic rotation equalizes the marginals. Lazy consistent $k$-coordinate sampling uses $O(k\log n)$ cached nodes after the LP is solved. Full output still costs $\Omega(n)$. For a noncommutative group this covers product identity, not an arbitrary fixed product: rotation conjugates that product. No succinct exponential graph or uncharged acquisition is hidden.

\paragraph{Exact precision obstruction.}
An integer straight-line program can be encoded by an explicit sparse rational quartic on a cube with Hessian between $\tfrac12I$ and $\tfrac32I$, a unique optimizer uniformly in the interior, and a coordinate sign encoding PosSLP. Repeated squaring produces doubly exponential small signed gaps and exponentially long rational optimizer representations. Approximation to a supplied number of bits remains polynomial; exact sign does not follow from contraction or conditioning alone. This is a scoped normal-form deduction, not a new complexity separation; broad PosSLP reductions to related fixed-point/convex interfaces already exist.

\paragraph{A tensor-compatible resolvent shortcut fails.}
OpenAI family 277 uses the resolvent Jensen remainder for $f(t)=t\log t$. Replacing its feature by the monoidal square root with coefficient one is false in operator order. Take
\[
C=\begin{pmatrix}1&0\\0&1/4\end{pmatrix},\qquad
P=\frac1{25}\begin{pmatrix}9&-12\\-12&16\end{pmatrix},\quad B_q=qP+(1-q)C.
\]
For the proposed Jensen gap minus square-root variance, $D(0)=0$ and, writing $\ell=\log4$,
\[
D'(0)=\begin{pmatrix}0&(4\ell-6)/25\\(4\ell-6)/25&16(\ell-1)/25\end{pmatrix},
\quad\det D'(0)=-(4\ell-6)^2/625<0.
\]
Thus $D(q)$ is not positive semidefinite for all sufficiently small positive $q$. This blocks that precise tensor-compatible shortcut, not all constants, representations, or multiplayer repetition theorems.

\paragraph{A prior-art stop and a physical scheduling stop.}
The all-characteristic selector extracted from family 198 was independently rederived, then found explicitly in the October 7 revision of \href{https://hexagonmath.org/2610.00140v1}{A note of refinements to an algebra of infinite little finitistic dimension}. It is retained as prior art, not a contribution. Treating a buffered multi-stage replica circuit as one channel improves a final-flag sample count, but does not provide early external flags or the LOCC causal interface required by the separate Astra gate. Neither route was promoted to a foundational result.

\section{What remains open and how to resume}

The completed central arguments contain no deliberately deferred mathematical lemma within their stated finite-dimensional, compact-source, or strong-curvature/score interfaces. Their external correctness and priority remain unestablished. The major-consequence standard is also unresolved: mixed-state classicalization, its sharp marginal separation, and transport stability are substantive, but this investigation has not established their historical originality or a transformative computational consequence with acquired inputs.

The hardest surviving routes are (i) acquisition-inclusive algorithms and the sharp universal product-cloning constant, (ii) quantitative dimension-dependent recovery under a marginal broadcasting promise, since a dimension-free modulus is disproved, (iii) source classes lacking both bounded support and strong curvature, or arbitrary radial Gaussian moment profiles beyond the completed power and logarithmic phase theorems, and (iv) approximate finite-semigroup coupling and the algebraic boundary excluding the exact Markov-cutpoint obstruction. Each requires new mathematics; none follows merely by changing names in the present proofs. The arbitrary Gaussian profile, weak-score boundary variant, and optional aperiodic-semigroup refinement are preserved as pending independent reconstruction. Precise terminal status is recorded in \texttt{research\_state.json}.

For continuation, begin with Theorem~\ref{thm:mixed}'s risk-gap induction, Theorem~\ref{thm:gaussian}'s quadratic-score interpolation, or the exact bridge (S2), then the corresponding companion proof. Keep pinned source status and imported premises separate. Do not restart a catalogue or silently upgrade a source claim because an internal derivation, finite check, or file hash succeeds.
\end{document}
